\documentclass[../../main.tex]{subfiles}

\graphicspath{{../../image/}{../../figure/}}

\begin{document}

\section{平衡律的表述}\label{section:ch01-平衡律的表述}

设 $\mathscr{X}$ 为 $\mathbb{R}^{k}$ 的一个开子集。$\mathscr{X}$ 中的一个\emph{真域}（proper domain）是指 $\mathscr{X}$ 的任意开的、有界的、具 Lipschitz 边界的子集。$\mathscr{X}$ 上的一条\emph{平衡律}假定：在任一真域 $\mathscr{D}$ 内某一“广延”量的\emph{生成}（production），由该量通过 $\mathscr{D}$ 的边界 $\partial\mathscr{D}$ 的\emph{通量}（flux）所平衡。

广延量的一个显著特征是：其生成与其通量在不相交子集上均可加。因此，真域 $\mathscr{D}$ 内的生成由 $\mathscr{X}$ 上的一个（带符号）Radon 测度 $\mathscr{P}$ 的值 $\mathscr{P}(\mathscr{D})$ 给出。类似地，对每个真域 $\mathscr{D}$，伴随一个定义在 $\partial\mathscr{D}$ 的 Borel 子集上的可数可加集函数 $\mathscr{Q}_{\mathscr{D}}$，使得通过 $\partial\mathscr{D}$ 的任一 Borel 子集 $\mathscr{C}$ 流入或流出 $\mathscr{D}$ 的通量由 $\mathscr{Q}_{\mathscr{D}}(\mathscr{C})$ 给出。于是平衡律简单地陈述为
\begin{align}
\mathscr{Q}_{\mathscr{D}}(\partial\mathscr{D})=\mathscr{P}(\mathscr{D}),\label{align:ch01-balance-primitive}
\end{align}
对 $\mathscr{X}$ 中每一真域 $\mathscr{D}$ 成立。

就本书的目的而言，只须考虑关于 Hausdorff 测度 $\mathscr{H}^{k-1}$ 绝对连续的通量集函数 $\mathscr{Q}_{\mathscr{D}}$。于是对 $\mathscr{X}$ 中每个真域 $\mathscr{D}$，伴随一个密度通量函数 $q_{\mathscr{D}}\in L^{1}(\partial\mathscr{D})$，使得
\begin{align}
\mathscr{Q}_{\mathscr{D}}(\mathscr{C})=\int_{\mathscr{C}}q_{\mathscr{D}}(X)\,\mathrm{d}\mathscr{H}^{k-1}(X),\label{align:ch01-flux-density}
\end{align}
对 $\partial\mathscr{D}$ 的任一 Borel 子集 $\mathscr{C}$ 成立。

$\partial\mathscr{D}$ 的 Borel 子集 $\mathscr{C}$ 通过 $\mathscr{C}$ 上各点处指向 $\mathscr{D}$ 外的单位法向量 $N$ 定向。平衡律理论中的基本公设是：通量仅依赖于曲面及其定向，即若 $\mathscr{C}$ 同时是两个不同真域 $\mathscr{D}_{1}$ 与 $\mathscr{D}_{2}$ 的边界的 Borel 子集，且二者在 $\mathscr{C}$ 上具有相同的外法向，则 $\mathscr{Q}_{\mathscr{D}_{1}}(\mathscr{C})=\mathscr{Q}_{\mathscr{D}_{2}}(\mathscr{C})$，从而对 $\mathscr{H}^{k-1}$ 几乎所有 $X\in\mathscr{C}$，有 $q_{\mathscr{D}_{1}}(X)=q_{\mathscr{D}_{2}}(X)$。

类比于通量测度，人们可能会想把考虑限制在关于 Lebesgue 测度绝对连续的生成测度 $\mathscr{P}$，从而由一个\emph{生成密度函数} $p\in L^{1}_{\mathrm{loc}}(\mathscr{X})$ 表示：
\begin{align}
\mathscr{P}(\mathscr{D})=\int_{\mathscr{D}}p(X)\,\mathrm{d}X.\label{align:ch01-production-density}
\end{align}

\end{document}
